\subsubsection{LeNet-5 Performance}\label{sec:lenet-5-performance}

LeNet-5 demonstrates that a CNN can successfully learn \hl{both feature extraction and classification directly from image data.} Its effectiveness on handwritten-digit recognition validates the key CNN idea: instead of manually designing image features, convolutional filters can be \textbf{learned end-to-end from data}.

\begin{figure}[!htp]
\centering
\tikzsetnextfilename{lenet-5-performance}
\begin{tikzpicture}[
    font=\sffamily,
    every node/.style={transform shape}
]

    % =========================================================================
    % 1. PARAMETERS & DIMENSIONS
    % =========================================================================
    \def\cellw{0.82}   % Cell width (cm)
    \def\cellh{0.58}   % Cell height (cm)
    \def\maxval{1128.0} % Maximum value for normalization

    % Main Title
    % \node[anchor=south west, font=\LARGE\bfseries, text=black] 
    %     at (-0.8, 0.6) {LeNet-5 performance};

    % =========================================================================
    % 2. CONFUSION MATRIX HEATMAP CELLS
    % =========================================================================
    \foreach \r/\c/\v in {
        0/0/975.0, 0/1/0.0, 0/2/0.0, 0/3/0.0, 0/4/0.0, 0/5/0.0, 0/6/2.0, 0/7/0.0, 0/8/2.0, 0/9/1.0,
        1/0/0.0, 1/1/1128.0, 1/2/2.0, 1/3/0.0, 1/4/0.0, 1/5/1.0, 1/6/1.0, 1/7/0.0, 1/8/3.0, 1/9/0.0,
        2/0/1.0, 2/1/1.0, 2/2/1014.0, 2/3/3.0, 2/4/1.0, 2/5/0.0, 2/6/1.0, 2/7/4.0, 2/8/7.0, 2/9/0.0,
        3/0/0.0, 3/1/0.0, 3/2/1.0, 3/3/1000.0, 3/4/0.0, 3/5/4.0, 3/6/0.0, 3/7/0.0, 3/8/4.0, 3/9/1.0,
        4/0/1.0, 4/1/0.0, 4/2/1.0, 4/3/0.0, 4/4/974.0, 4/5/0.0, 4/6/1.0, 4/7/0.0, 4/8/2.0, 4/9/3.0,
        5/0/2.0, 5/1/0.0, 5/2/0.0, 5/3/5.0, 5/4/0.0, 5/5/882.0, 5/6/1.0, 5/7/0.0, 5/8/1.0, 5/9/1.0,
        6/0/4.0, 6/1/2.0, 6/2/0.0, 6/3/0.0, 6/4/4.0, 6/5/5.0, 6/6/942.0, 6/7/0.0, 6/8/1.0, 6/9/0.0,
        7/0/0.0, 7/1/4.0, 7/2/4.0, 7/3/1.0, 7/4/2.0, 7/5/0.0, 7/6/0.0, 7/7/1004.0, 7/8/1.0, 7/9/12.0,
        8/0/3.0, 8/1/0.0, 8/2/1.0, 8/3/0.0, 8/4/1.0, 8/5/1.0, 8/6/0.0, 8/7/2.0, 8/8/962.0, 8/9/4.0,
        9/0/1.0, 9/1/2.0, 9/2/0.0, 9/3/0.0, 9/4/7.0, 9/5/4.0, 9/6/0.0, 9/7/3.0, 9/8/1.0, 9/9/991.0
    }{
        % Normalize color intensity (0% to 100%)
        \pgfmathsetmacro{\fperc}{(\v > 0 ? ((\v / \maxval) * 100) : 0)}
        
        % Invert text color to white for dark cells
        \pgfmathparse{\v > 400 ? 1 : 0}
        \ifnum\pgfmathresult=1
            \def\txtcol{white}
        \else
            \def\txtcol{black}
        \fi

        % Draw cell background & border
        \filldraw[draw=gray!55, line width=0.4pt, fill=deepgreen!\fperc!white]
            (\c*\cellw, -\r*\cellh) rectangle ++(\cellw, -\cellh);

        % Display numeric value
        \node[font=\fontsize{6.5pt}{8pt}\selectfont, text=\txtcol] 
            at (\c*\cellw + 0.5*\cellw, -\r*\cellh - 0.5*\cellh) {\v};
    }

    % =========================================================================
    % 3. AXIS LABELS & TICKS
    % =========================================================================
    
    % Y-Axis Ticks (True Classes: 0 to 9)
    \foreach \r in {0,...,9} {
        \node[font=\footnotesize, left=2pt] 
            at (0, -\r*\cellh - 0.5*\cellh) {\r};
    }
    % Y-Axis Title
    \node[rotate=90, font=\large\bfseries, anchor=center] 
        at (-0.8, -5*\cellh) {True classes};

    % X-Axis Ticks (Predicted Classes: 0 to 9)
    \foreach \c in {0,...,9} {
        \node[font=\footnotesize, below=2pt] 
            at (\c*\cellw + 0.5*\cellw, -10*\cellh) {\c};
    }
    % X-Axis Title
    \node[font=\large\bfseries, anchor=center] 
        at (5*\cellw, -10*\cellh - 0.65) {Predicted classes};

    % =========================================================================
    % 4. COLORBAR (RIGHT SIDE)
    % =========================================================================
    \def\cbx{10*\cellw + 0.5} % Colorbar x-position
    \def\cbw{0.28}            % Colorbar width
    \def\cbh{10*\cellh}       % Colorbar height

    % Gradient Colorbar Bar
    \shade[top color=deepgreen, bottom color=white, draw=none] 
        (\cbx, -\cbh) rectangle ++(\cbw, \cbh);

    % Colorbar Ticks and Labels (0 to 1000)
    \foreach \val in {0, 200, 400, 600, 800, 1000} {
        \pgfmathsetmacro{\ypos}{-\cbh + (\val / \maxval) * \cbh}
        \draw[black!70, line width=0.5pt] 
            (\cbx + \cbw, \ypos) -- ++(0.08, 0) 
            node[right=1pt, font=\scriptsize] {\val};
    }

\end{tikzpicture}
\caption{\textbf{LeNet-5 classification performance.} Confusion matrix on the $10'000$ test images, where rows indicate true classes and columns predicted classes. The strongly dominant diagonal shows that most digits are correctly classified, with $9'872$ correct predictions and only $128$ errors, corresponding to an accuracy of 98.72\%. The few off-diagonal entries indicate occasional confusion between visually similar handwritten digits.}
\label{fig:lenet-5-performance}
\end{figure}

\noindent
The \autoref{fig:lenet-5-performance} provides experimental evidence for the entire CNN idea developed throughout the chapter:
\begin{equation*}
    \text{image}
    \rightarrow
    \text{learned local features}
    \rightarrow
    \text{higher-level representation}
    \rightarrow
    \text{classification}
\end{equation*}
LeNet does not need manually designed digit features. The convolutional feature extractor learns useful representations directly from the images, and the resulting network achieves \textbf{98.72\% test accuracy} on the dataset represented by this confusion matrix.